% GENERATED FILE. Edit canonical lessons, then run npm run generate:books.
% canonical-source-hash: fnv1a64:d321fc64d70e5102
% canonical-lessons: PA-C330-bharjai, PA-C330-jija, PA-C330-nanan, PA-C330-darani, PA-C330-jathani

\chapter{Brothers' and Sisters' Spouses}
\label{ch:pa-a2-brothers-and-sisters-spouses}
\hlchaptermodality{330}

\begin{chapteropening}
\textbf{By the end of this chapter:} \emph{I can name the husbands and wives of my brothers and sisters.}

Complete the last lesson of chapter 330: I can name the husbands and wives of my brothers and sisters.
\end{chapteropening}

\section[\texorpdfstring{\pa{ਭਰਜਾਈ} (bharjāī)}{ਭਰਜਾਈ (bharjāī)}]{\pa{ਭਰਜਾਈ} (bharjāī) — a brother's wife}

\label{lesson:PA-C330-bharjai}



\begin{quote}
Before the new one: say the Punjabi for a son-in-law, then the Punjabi for a child's father-in-law.
\end{quote}

\subsection*{You'll want to know: \pa{ਭਰਜਾਈ}}
\textbf{\pa{ਭਰਜਾਈ}} — \emph{bharjāī} — \textquotedblleft{}a brother's wife\textquotedblright{}.

\begin{quote}
\pa{ਮੇਰੀ} \pa{ਭਰਜਾਈ} \pa{ਸਕੂਲ} \pa{ਵਿੱਚ} \pa{ਬੱਚਿਆਂ} \pa{ਨੂੰ} \pa{ਪੰਜਾਬੀ} \pa{ਸਿਖਾਉਂਦੀ} \pa{ਹੈ।}
\end{quote}

\begin{quote}
\emph{merī bharjāī sakūl vich bachchiāṁ nū̃ panjābī sikhāundī hai.}
\end{quote}

\begin{quote}
My brother's wife teaches Punjabi to children at school.
\end{quote}

A feminine noun.

\begin{cousinweb}
\textbf{\pa{ਭਰਜਾਈ}} (\emph{bharjāī}) is inherited from Sanskrit \emph{bhrāturjāyā}, \textquotedblleft{}a brother's wife\textquotedblright{}; CDIAL lists Punjabi \emph{bharjāī}, \textquotedblleft{}brother's wife\textquotedblright{}, under it (CDIAL 9660).
\end{cousinweb}

\begin{grammarlens}[title={What you've built}]
One more word for this chapter.
\end{grammarlens}

\subsection*{Your turn}
\begin{itemize}
\raggedright
  \item \emph{Say it:} \emph{bharjāī}
  \item \emph{Say it:} \emph{bharjāī}, once more
  \item \emph{Say it:} \emph{merī bharjāī sakūl vich bachchiāṁ nū̃ panjābī sikhāundī hai.} — My brother's wife teaches Punjabi to children at school.
  \item \emph{Recall:} say the Punjabi for a son-in-law, then the Punjabi for a child's father-in-law, then say \emph{bharjāī} again
\end{itemize}

\begin{tcolorbox}[breakable,colback=teal!4,colframe=teal!35!black,title={Before you move on}]
What does \pa{ਭਰਜਾਈ} mean? (\textquotedblleft{}A brother's wife\textquotedblright{}.) Say it once more.
\end{tcolorbox}
\section[\texorpdfstring{\pa{ਜੀਜਾ} (jījā)}{ਜੀਜਾ (jījā)}]{\pa{ਜੀਜਾ} (jījā) — a sister's husband}

\label{lesson:PA-C330-jija}



\begin{quote}
Before the new one: say the Punjabi for a child's father-in-law, then the Punjabi for a brother's wife.
\end{quote}

\subsection*{You'll want to know: \pa{ਜੀਜਾ}}
\textbf{\pa{ਜੀਜਾ}} — \emph{jījā} — \textquotedblleft{}a sister's husband\textquotedblright{}.

\begin{quote}
\pa{ਮੇਰਾ} \pa{ਜੀਜਾ} \pa{ਬਹੁਤ} \pa{ਹੱਸਦਾ} \pa{ਹੈ।}
\end{quote}

\begin{quote}
\emph{merā jījā bahut hassdā hai.}
\end{quote}

\begin{quote}
My sister's husband laughs a lot.
\end{quote}

A masculine noun.

\begin{cousinweb}
\textbf{\pa{ਜੀਜਾ}} (\emph{jījā}) is inherited from a reconstructed \emph{*jījja}, \textquotedblleft{}a nursery word for breast, mother, sister, relative\textquotedblright{}; CDIAL lists Punjabi \emph{jījā}, \textquotedblleft{}sister's husband\textquotedblright{}, under it (CDIAL 5232).
\end{cousinweb}

\begin{grammarlens}[title={What you've built}]
One more word for this chapter.
\end{grammarlens}

\subsection*{Your turn}
\begin{itemize}
\raggedright
  \item \emph{Say it:} \emph{jījā}
  \item \emph{Say it:} \emph{jījā}, once more
  \item \emph{Say it:} \emph{merā jījā bahut hassdā hai.} — My sister's husband laughs a lot.
  \item \emph{Recall:} say the Punjabi for a child's father-in-law, then the Punjabi for a brother's wife, then say \emph{jījā} again
  \item \emph{Write it:} copy \textbf{\pa{ਜੀਜਾ}} once, checking the whole word
\end{itemize}

\begin{tcolorbox}[breakable,colback=teal!4,colframe=teal!35!black,title={Before you move on}]
What does \pa{ਜੀਜਾ} mean? (\textquotedblleft{}A sister's husband\textquotedblright{}.) Say it once more.
\end{tcolorbox}
\section[\texorpdfstring{\pa{ਨਣਾਨ} (naṇān)}{ਨਣਾਨ (naṇān)}]{\pa{ਨਣਾਨ} (naṇān) — a husband's sister}

\label{lesson:PA-C330-nanan}



\begin{quote}
Before the new one: say the Punjabi for a brother's wife, then the Punjabi for a sister's husband.
\end{quote}

\subsection*{You'll want to know: \pa{ਨਣਾਨ}}
\textbf{\pa{ਨਣਾਨ}} — \emph{naṇān} — \textquotedblleft{}a husband's sister\textquotedblright{}.

\begin{quote}
\pa{ਮੇਰੀ} \pa{ਨਣਾਨ} \pa{ਹਰ} \pa{ਐਤਵਾਰ} \pa{ਸਾਡੇ} \pa{ਘਰ} \pa{ਆਉਂਦੀ} \pa{ਹੈ।}
\end{quote}

\begin{quote}
\emph{merī naṇān har aitvār sāḍe ghar āundī hai.}
\end{quote}

\begin{quote}
My husband's sister comes to our house every Sunday.
\end{quote}

A feminine noun.

\begin{cousinweb}
\textbf{\pa{ਨਣਾਨ}} (\emph{naṇān}) is inherited from Sanskrit \emph{nanāndṛ}, \textquotedblleft{}a husband's sister\textquotedblright{}; CDIAL lists Punjabi \emph{nanāṇ}, \textquotedblleft{}husband's sister\textquotedblright{}, under it (CDIAL 6946).
\end{cousinweb}

\begin{grammarlens}[title={What you've built}]
One more word for this chapter.
\end{grammarlens}

\subsection*{Your turn}
\begin{itemize}
\raggedright
  \item \emph{Say it:} \emph{naṇān}
  \item \emph{Say it:} \emph{naṇān}, once more
  \item \emph{Say it:} \emph{merī naṇān har aitvār sāḍe ghar āundī hai.} — My husband's sister comes to our house every Sunday.
  \item \emph{Recall:} say the Punjabi for a brother's wife, then the Punjabi for a sister's husband, then say \emph{naṇān} again
  \item \emph{Write it:} copy \textbf{\pa{ਨਣਾਨ}} once, checking the whole word
\end{itemize}

\begin{tcolorbox}[breakable,colback=teal!4,colframe=teal!35!black,title={Before you move on}]
What does \pa{ਨਣਾਨ} mean? (\textquotedblleft{}A husband's sister\textquotedblright{}.) Say it once more.
\end{tcolorbox}
\section[\texorpdfstring{\pa{ਦਰਾਣੀ} (darāṇī)}{ਦਰਾਣੀ (darāṇī)}]{\pa{ਦਰਾਣੀ} (darāṇī) — a husband's younger brother's wife}

\label{lesson:PA-C330-darani}



\begin{quote}
Before the new one: say the Punjabi for a sister's husband, then the Punjabi for a husband's sister.
\end{quote}

\subsection*{You'll want to know: \pa{ਦਰਾਣੀ}}
\textbf{\pa{ਦਰਾਣੀ}} — \emph{darāṇī} — \textquotedblleft{}a husband's younger brother's wife\textquotedblright{}.

\begin{quote}
\pa{ਮੇਰੀ} \pa{ਦਰਾਣੀ} \pa{ਮੇਰੇ} \pa{ਨਾਲ} \pa{ਰਸੋਈ} \pa{ਵਿੱਚ} \pa{ਕੰਮ} \pa{ਕਰਦੀ} \pa{ਹੈ।}
\end{quote}

\begin{quote}
\emph{merī darāṇī mere nāl rasoī vich kamm kardī hai.}
\end{quote}

\begin{quote}
My husband's younger brother's wife works with me in the kitchen.
\end{quote}

A feminine noun.

\begin{cousinweb}
\textbf{\pa{ਦਰਾਣੀ}} (\emph{darāṇī}) is inherited from a reconstructed \emph{*devarajāni}, \textquotedblleft{}a husband's younger brother's wife\textquotedblright{}; CDIAL lists Punjabi \emph{darāṇī}, \textquotedblleft{}husband's younger brother's wife\textquotedblright{}, under it (CDIAL 6534).
\end{cousinweb}

\begin{grammarlens}[title={What you've built}]
One more word for this chapter.
\end{grammarlens}

\subsection*{Your turn}
\begin{itemize}
\raggedright
  \item \emph{Say it:} \emph{darāṇī}
  \item \emph{Say it:} \emph{darāṇī}, once more
  \item \emph{Say it:} \emph{merī darāṇī mere nāl rasoī vich kamm kardī hai.} — My husband's younger brother's wife works with me in the kitchen.
  \item \emph{Recall:} say the Punjabi for a sister's husband, then the Punjabi for a husband's sister, then say \emph{darāṇī} again
  \item \emph{Write it:} copy \textbf{\pa{ਦਰਾਣੀ}} once, checking the whole word
\end{itemize}

\begin{tcolorbox}[breakable,colback=teal!4,colframe=teal!35!black,title={Before you move on}]
What does \pa{ਦਰਾਣੀ} mean? (\textquotedblleft{}A husband's younger brother's wife\textquotedblright{}.) Say it once more.
\end{tcolorbox}
\section[\texorpdfstring{\pa{ਜਠਾਣੀ} (jaṭhāṇī)}{ਜਠਾਣੀ (jaṭhāṇī)}]{\pa{ਜਠਾਣੀ} (jaṭhāṇī) — a husband's elder brother's wife}

\label{lesson:PA-C330-jathani}



\begin{quote}
Before the new one: say the Punjabi for a husband's sister, then the Punjabi for a husband's younger brother's wife.
\end{quote}

\subsection*{You'll want to know: \pa{ਜਠਾਣੀ}}
\textbf{\pa{ਜਠਾਣੀ}} — \emph{jaṭhāṇī} — \textquotedblleft{}a husband's elder brother's wife\textquotedblright{}.

\begin{quote}
\pa{ਮੇਰੀ} \pa{ਜਠਾਣੀ} \pa{ਬਹੁਤ} \pa{ਸਿਆਣੀ} \pa{ਹੈ।}
\end{quote}

\begin{quote}
\emph{merī jaṭhāṇī bahut siāṇī hai.}
\end{quote}

\begin{quote}
My husband's elder brother's wife is very wise.
\end{quote}

A feminine noun.

\begin{cousinweb}
\textbf{\pa{ਜਠਾਣੀ}} (\emph{jaṭhāṇī}) is inherited from a reconstructed \emph{*jyeṣṭhajāni}, \textquotedblleft{}a senior wife; an elder brother's wife\textquotedblright{}; CDIAL lists Punjabi \emph{jaṭhāṇī}, \textquotedblleft{}husband's elder brother's wife\textquotedblright{}, under it (CDIAL 5287).
\end{cousinweb}

\begin{grammarlens}[title={What you've built}]
One more word for this chapter.
\end{grammarlens}

\subsection*{Your turn}
\begin{itemize}
\raggedright
  \item \emph{Say it:} \emph{jaṭhāṇī}
  \item \emph{Say it:} \emph{jaṭhāṇī}, once more
  \item \emph{Say it:} \emph{merī jaṭhāṇī bahut siāṇī hai.} — My husband's elder brother's wife is very wise.
  \item \emph{Recall:} say the Punjabi for a husband's sister, then the Punjabi for a husband's younger brother's wife, then say \emph{jaṭhāṇī} again
  \item \emph{Write it:} copy \textbf{\pa{ਜਠਾਣੀ}} once, checking the whole word
\end{itemize}

\begin{tcolorbox}[breakable,colback=teal!4,colframe=teal!35!black,title={Before you move on}]
What does \pa{ਜਠਾਣੀ} mean? (\textquotedblleft{}A husband's elder brother's wife\textquotedblright{}.) Say it once more.
\end{tcolorbox}
